% =============================================================================
% CHAPTER 5 — RESEARCH METHODOLOGY
% Thesis: PATHCAST
% Version: REDUCED — paragraph-level paraphrase of tables removed.
%          All formal content (definitions, equations, protocols) preserved.
%          Cross-references unchanged.
% =============================================================================

\chapter{Research Methodology}
\label{ch:methodology}

This chapter describes the methodological framework. Section~\ref{sec:dsr}
presents Design Science Research. Section~\ref{sec:dataset} describes the
dataset. Section~\ref{sec:eval-design} formalizes the experimental design.
Section~\ref{sec:evaluation-metrics} defines evaluation metrics.
Section~\ref{sec:baselines} specifies baselines.
Section~\ref{sec:statistical-analysis} describes statistical procedures.
Section~\ref{sec:threats} discusses threats to validity.


% =============================================================================
% 5.1 DESIGN SCIENCE RESEARCH
% =============================================================================
\section{Design Science Research Paradigm}
\label{sec:dsr}

This research adopts Design Science Research (DSR) following
Hevner et al.~\cite{hevner2004} and the DSRM process of Peffers et
al.~\cite{peffers2007}.

\subsection{DSR Artifact Characterization}
\label{sec:dsr-artifact}

The artifact produced is a \textbf{method}: PATHCAST is a formally
specified pipeline that prescribes how to transform repository data into
probabilistic process forecasts. The supporting \textbf{instantiation} is
a Python tool implementing the four pipeline stages and the ML component.


\subsection{DSRM Process Instantiation}
\label{sec:dsrm-instantiation}

\begin{table}[htbp]
\centering
\caption{Mapping of DSRM activities to thesis components.}
\label{tab:dsrm-mapping}
\small
\begin{tabular}{clll}
\toprule
\textbf{Activity} & \textbf{DSRM Activity} & \textbf{Chapter} & \textbf{Deliverable} \\
\midrule
1 & Problem Identification & Ch.~\ref{ch:introduction} & Three research problems \\
  &                        & Ch.~\ref{ch:theoretical-foundation} & Gap analysis table \\
\midrule
2 & Objectives Definition  & Ch.~\ref{ch:introduction} & 3 RQs, 4 objectives, 5 contributions \\
\midrule
3 & Design \& Development  & Ch.~\ref{ch:method}        & PATHCAST method \\
  &                        & Ch.~\ref{ch:implementation} & Python tool \\
\midrule
4 & Demonstration          & Ch.~\ref{ch:evaluation}    & 48 repositories \\
\midrule
5 & Evaluation             & Ch.~\ref{ch:evaluation}    & E1: Pipeline; E2: Forecasting \\
\midrule
6 & Communication          & This thesis                & Conference/journal papers \\
\bottomrule
\end{tabular}
\end{table}


\subsection{DSR Knowledge Contribution}
\label{sec:dsr-contribution}

Following Gregor and Hevner~\cite{gregor2013positioning}, this research
operates in the \emph{Improvement} quadrant: a known problem class
(software process forecasting) addressed with a new solution
(process-aware stochastic simulation parameterized by process mining).
The contribution is the formal integration of established techniques
into a SDLC-specific pipeline, not the invention of new individual
techniques.


% =============================================================================
% 5.2 DATASET
% =============================================================================
\section{Dataset Description}
\label{sec:dataset}

The empirical evaluation uses 48 open-source software repositories.

\subsection{Repository Selection Criteria}
\label{sec:selection-criteria}

Inclusion criteria:
\begin{enumerate}
    \item \textbf{Minimum commit volume:} $\geq 1{,}000$ commits in main.
    \item \textbf{Issue tracker usage:} GitHub Issues or Jira with
    $\geq 100$ closed issues containing state transitions.
    \item \textbf{Pull request workflow:} evidence of PR-based development.
    \item \textbf{Activity period:} $\geq 12$ months continuous activity.
    \item \textbf{Non-trivial contributors:} $\geq 5$ distinct contributors.
    \item \textbf{Non-fork:} not a fork of another project.
\end{enumerate}

Exclusion: $>50\%$ bot commits; external-bug-only issue tracking; mirror
or archive repositories.


\subsection{Selection Process}
\label{sec:selection-process}

\paragraph{Step 1 — Candidate generation.} 150 candidates from SEART
GitHub Search~\cite{seart2021}: $\geq 1{,}000$ commits, issues+PRs,
non-fork, last push $\geq$~2022-01-01, languages
\{Python, JS/TS, Java, Go, Rust, C/C++\}.

\paragraph{Step 2 — Criterion verification.} Each candidate verified
against C2 and C5 via the GitHub REST API~\cite{github2023api}.

\paragraph{Step 3 — Stratified selection.} 48 repositories selected by
stratified sampling across the dimensions of
Table~\ref{tab:dataset-diversity}: 3--4 per primary language; balanced
across scale tiers; $\geq 4$ per governance model.

\paragraph{Step 4 — Jira enrichment.} For Apache and Spring projects,
the issue component was enriched with the Public Jira Dataset of
Montgomery et al.~\cite{montgomery2022jira}.

The complete list of 48 repositories is in
Appendix~\ref{app:repository-list}.


\subsection{Data Collection}
\label{sec:data-collection}

Per repository:
\begin{itemize}
    \item Git commit history via
    \texttt{/repos/\{owner\}/\{repo\}/commits}.
    \item Issue lifecycle via the timeline events API
    (or Jira~\cite{montgomery2022jira} for Apache/Spring).
    \item Pull requests via GraphQL API.
    \item GitHub Actions runs (only $\geq 50$ runs counts as Advanced
    maturity).
\end{itemize}

Data stored in PostgreSQL for Stage~1 ingestion.


\subsection{Dataset Descriptive Statistics}
\label{sec:dataset-stats}

\begin{table}[htbp]
\centering
\caption{Aggregate dataset statistics (48 repositories).}
\label{tab:dataset-summary}
\begin{tabular}{lrrrr}
\toprule
\textbf{Metric} & \textbf{Total} & \textbf{Mean} & \textbf{Median} & \textbf{Std Dev} \\
\midrule
Commits                 & 227{,}000 & 4{,}729 & 3{,}215 & -- \\
Closed issues           & --        & --      & --      & -- \\
Pull requests           & --        & --      & --      & -- \\
Contributors            & --        & --      & --      & -- \\
Activity span (months)  & --        & --      & --      & -- \\
\bottomrule
\multicolumn{5}{l}{\parbox{13cm}{\footnotesize ``--'' populated upon completion of
    data collection. 44 of 48 repositories confirmed at writing.}}
\end{tabular}
\end{table}

\begin{table}[htbp]
\centering
\caption{Dataset diversity profile.}
\label{tab:dataset-diversity}
\begin{tabular}{ll}
\toprule
\textbf{Dimension} & \textbf{Coverage} \\
\midrule
Primary languages   & Python, JS/TS, Java, Go, Rust, C/C++ \\
Domain categories   & Web, CLI, libraries, DevOps, data processing \\
Project scale       & Small ($<5$K), Medium (5--20K), Large ($>20$K commits) \\
Governance model    & Foundation, corporate, community-driven \\
Workflow maturity   & Basic / Standard / Advanced \\
Issue source        & GitHub Issues; Jira via~\cite{montgomery2022jira} \\
\bottomrule
\end{tabular}
\end{table}


\subsection{Ethical Considerations}
\label{sec:ethics}

All data is publicly available under open-source licenses. No private
repositories, proprietary data, or personally identifiable information
beyond public GitHub usernames. Compliant with GitHub Terms of Service.


% =============================================================================
% 5.3 EXPERIMENTAL DESIGN
% =============================================================================
\section{Experimental Design}
\label{sec:eval-design}

\begin{figure}[htbp]
\centering
\begin{tikzpicture}[
    node distance=1.5cm,
    stage/.style={rectangle, draw, thick, rounded corners=4pt,
                  minimum width=5cm, minimum height=1.4cm,
                  align=center, font=\small},
    rq/.style={rectangle, draw, dashed, rounded corners=2pt,
               minimum width=2.2cm, minimum height=0.8cm,
               align=center, font=\scriptsize, fill=gray!5},
    arrow/.style={->, thick, >=stealth}
]
    \node[stage, fill=blue!10] (e1) {
        \textbf{E1: Pipeline Validation}\\
        {\scriptsize Event log + Process model quality}
    };
    \node[rq, right=1cm of e1] (rq1) {RQ1};
    \draw[arrow] (e1) -- (rq1);

    \node[stage, fill=orange!10, below=1.5cm of e1] (e2) {
        \textbf{E2: Forecasting Accuracy}\\
        {\scriptsize Stochastic forecasts + ML comparison}
    };
    \node[rq, right=1cm of e2, yshift=0.4cm] (rq2) {RQ2};
    \node[rq, right=1cm of e2, yshift=-0.4cm] (rq3) {RQ3};
    \draw[arrow] (e2) -- (rq2);
    \draw[arrow] (e2) -- (rq3);

    \draw[arrow] (e1) -- (e2) node[midway, right, font=\scriptsize]
        {Validated event logs};

    \node[rectangle, draw, fill=green!8, rounded corners=2pt,
          minimum width=3cm, font=\scriptsize, align=center,
          left=1.5cm of e1, yshift=-0.75cm] (data) {
        48 Repositories\\227K commits\\Dataset $\mathcal{R}$
    };
    \draw[arrow] (data) |- (e1);
    \draw[arrow] (data) |- (e2);
\end{tikzpicture}
\caption{Two-stage experimental design.}
\label{fig:eval-design}
\end{figure}


\subsection{E1: Pipeline Validation}
\label{sec:e1-design}

E1 assesses whether Stages 1 and 2 produce event logs and process models
of sufficient quality. Addresses RQ1.

\subsubsection{Scope}

\begin{enumerate}
    \item Event log quality: correlation coverage, activity coverage,
    trace completeness.
    \item Process model quality: fitness, precision, generalization via
    alignment-based conformance (Section~\ref{sec:s2-conformance}).
\end{enumerate}

\subsubsection{Protocol}

For each $r_i \in \mathcal{R}$:
\begin{enumerate}
    \item Execute Stage~1 with the default correlation configuration.
    \item Compute event log quality metrics
    (Definition~\ref{def:log-quality}).
    \item Execute Stage~2 with IMf at $\theta_{\text{noise}} = 0.2$.
    \item Compute fitness, precision, generalization.
    \item Record process variants and state space $|S|$.
    \item Flag transitions with $< 30$ observations; apply Laplace
    smoothing ($\alpha = 1$) at Stage~3.
\end{enumerate}

\subsubsection{Sensitivity Analysis}

$\theta_{\text{noise}}$ varied over $\{0.0, 0.1, 0.2, 0.3, 0.5\}$ on 10
representative repositories. Pareto front of fitness vs.\ precision
identifies the threshold balancing both objectives.

\subsubsection{Acceptance Criteria}

\begin{equation}
\label{eq:acceptance}
\text{fitness}(\mathcal{L}, M) \geq 0.7 \;\wedge\;
\text{coverage}(\mathcal{L}) \geq 0.5 \;\wedge\;
|S_T| \geq 2 \;\wedge\;
|\mathcal{L}_{\text{test}}| \geq 50.
\end{equation}

The fourth condition reflects the variance threshold for empirical
percentiles and CRPS~\cite{gneiting2007strictly}. Repositories satisfying
the first three but failing the case count threshold are reported as
\emph{quality-validated but statistically insufficient for E2} and
excluded from E2 only.


\subsection{E2: Forecasting Accuracy}
\label{sec:e2-design}

E2 assesses Stages 3 and 4 forecasts and the predictive contribution of
process-derived ML features. Addresses RQ2 and RQ3.

\subsubsection{Scope}

For each repository passing E1:
\begin{enumerate}
    \item Execute Stages 3 and 4 to produce probabilistic lead time
    forecasts.
    \item Execute the ML component with three feature sets (PM-only,
    MSR-only, Combined).
    \item Compare against ground truth.
\end{enumerate}

\subsubsection{Temporal Hold-Out Protocol}
\label{sec:temporal-holdout}

\begin{definition}[Temporal Hold-Out]
\label{def:temporal-holdout}
For each repository with event log $\mathcal{L}$ over
$[t_{\min}, t_{\max}]$:
\begin{align}
t_{\text{cut}} &= t_{\min} + 0.7 \cdot (t_{\max} - t_{\min})
    \label{eq:cutoff} \\
\mathcal{L}_{\text{train}} &= \{c : \max(\tau(\sigma_c)) \leq t_{\text{cut}}\}
    \label{eq:train-set} \\
\mathcal{L}_{\text{test}}  &= \{c : \min(\tau(\sigma_c)) > t_{\text{cut}}\}
    \label{eq:test-set}
\end{align}
\end{definition}

The 70/30 split follows
convention~\cite{tashman2000out}; sensitivity to 60/40 and 80/20 is
assessed on a representative subset.

\textbf{Model estimation:} $P$, $\hat{F}$, ML models estimated from
$\mathcal{L}_{\text{train}}$ only.

\textbf{Forecast evaluation:} for each $c \in \mathcal{L}_{\text{test}}$,
forecast vs.\ observed $T_c^{\text{obs}}$, $y_c^{\text{obs}}$.


\subsubsection{Cross-Repository Generalization}
\label{sec:cross-repo}

Leave-one-project-out (LOPO) within domain groups (e.g., all Python web
frameworks). Minimum group size~5; smaller groups reported descriptively
only. Current dataset: Python web frameworks
(Django, Flask, FastAPI, Requests, Click) and Go DevOps tools
(kubectl, Helm, Terraform, Prometheus, cli) qualify; others descriptive.


% =============================================================================
% 5.4 EVALUATION METRICS
% =============================================================================
\section{Evaluation Metrics}
\label{sec:evaluation-metrics}

% --- RQ1 metrics ---
\subsection{Metrics for RQ1: Event Log and Process Model Quality}
\label{sec:metrics-rq1}

Event log quality and process model quality are reported separately to
isolate sparse-data effects from modeling effects.

\begin{table}[ht]
\centering
\caption{Evaluation metrics for RQ1.}
\label{tab:metrics-rq1}
\small
\begin{tabular}{@{}l l p{7.8cm}@{}}
\toprule
\textbf{Metric} & \textbf{Range} & \textbf{Definition} \\
\midrule
\multicolumn{3}{@{}l}{\textit{Event Log Quality}} \\[2pt]
Correlation Coverage & $[0,1]$ & Fraction of raw events assigned to a case (Eq.~4.15). \\
Activity Coverage    & $[0,1]$ & Fraction of SDLC taxonomy observed (Eq.~4.22). \\
Trace Completeness   & $[0,1]$ & Fraction of cases reaching absorbing state (Eq.~4.20). \\
Mean Trace Length    & $\mathbb{R}^+$ & Average events per trace (Eq.~4.21). \\
Variant Count        & $\mathbb{N}$ & Distinct activity sequences (Eq.~4.23). \\
\midrule
\multicolumn{3}{@{}l}{\textit{Process Model Quality}} \\[2pt]
Fitness              & $[0,1]$ & Alignment-based fitness (Eq.~2.2). Higher better. \\
Precision            & $[0,1]$ & Fraction of model behaviour observed in log. Higher better. \\
Generalization       & $[0,1]$ & Capacity to capture unseen legitimate behaviour. \\
\bottomrule
\end{tabular}
\end{table}

\paragraph{Metric rationale.} Correlation Coverage signals Stage~1
fidelity. Activity Coverage contextualizes fitness: high fitness on a
log with low activity coverage may reflect an oversimplified model.
Trace Completeness bounds the supervised-learning training set. Mean
Trace Length and Variant Count characterize structural complexity.

The four-quality framework~\cite{aalst2016process} additionally includes
\emph{simplicity}; this is constrained by construction in PATHCAST since
the state space is bounded by the SDLC taxonomy
(Section~\ref{sec:s1-mapping}) and not reported separately.


% --- RQ2 metrics ---
\subsection{Metrics for RQ2: Stochastic Forecasting Accuracy}
\label{sec:metrics-rq2}

\begin{table}[ht]
\centering
\caption{Evaluation metrics for RQ2.}
\label{tab:metrics-rq2}
\small
\begin{tabular}{@{}l l p{7.2cm}@{}}
\toprule
\textbf{Metric} & \textbf{Type} & \textbf{Definition} \\
\midrule
\multicolumn{3}{@{}l}{\textit{Point Forecast Accuracy}} \\[2pt]
MAE  & Error & $\tfrac{1}{n}\sum_c |T_c^{\text{obs}} - \hat{T}_c^{\text{P50}}|$ \\
MAPE & Error (\%) & $\tfrac{100}{n}\sum_c |T_c^{\text{obs}} - \hat{T}_c^{\text{P50}}| / T_c^{\text{obs}}$ \\
MdAE & Error & $\operatorname{median}(|T_c^{\text{obs}} - \hat{T}_c^{\text{P50}}|)$ \\
\midrule
\multicolumn{3}{@{}l}{\textit{Distributional Accuracy}} \\[2pt]
CRPS & Score & $\tfrac{1}{n}\sum_c \int [\hat{F}_c(t) - \mathbf{1}(t \geq T_c^{\text{obs}})]^2 dt$ \\
ACE  & Error & $\tfrac{1}{9}\sum_{\alpha \in \{0.1,\ldots,0.9\}} |\hat{p}_\alpha - \alpha|$ \\
Coverage 85\% & $[0,1]$ & Fraction of $T_c^{\text{obs}} \in [\hat{Q}_{0.075}, \hat{Q}_{0.925}]$. \\
\midrule
\multicolumn{3}{@{}l}{\textit{Outcome Prediction}} \\[2pt]
Completion Accuracy & $[0,1]$ & Done vs.\ Cancelled prediction accuracy. \\
\bottomrule
\end{tabular}
\end{table}

\paragraph{CRPS} is the primary metric. Reduces to MAE for point
forecasts~\cite{gneiting2007strictly}, preserving comparability with
deterministic baselines. Derivation in Appendix~B.

\paragraph{Point metrics.} MAE in calendar days; MAPE for cross-repository
comparison; MdAE robust to outliers.

\paragraph{ACE} averages absolute deviation across nine deciles. ACE = 0
denotes perfect calibration; ACE $> 0.1$ indicates systematic over- or
underconfidence. Reported alongside CRPS because a forecast can have low
CRPS while being miscalibrated.

\paragraph{Coverage at 85\%} is the standard agile commitment
interval~\cite{vacanti2015}; complements ACE by focusing on the single
operationally-relevant interval.

\paragraph{Completion Accuracy.} Predicted Done vs.\ Cancelled per
Definition~\ref{def:label}.


% --- RQ3 metrics ---
\subsection{Metrics for RQ3: ML Feature Comparison}
\label{sec:metrics-rq3}

\begin{table}[ht]
\centering
\caption{Evaluation metrics for RQ3.}
\label{tab:metrics-rq3}
\small
\begin{tabular}{@{}l l p{7.4cm}@{}}
\toprule
\textbf{Metric} & \textbf{Type} & \textbf{Definition} \\
\midrule
AUC-ROC & Discrimination & Threshold-independent discrimination. \\
F1-Score & Balance & $2 \cdot \mathrm{Precision} \cdot \mathrm{Recall} / (\mathrm{Precision} + \mathrm{Recall})$ at validation-optimal threshold. \\
Brier Score & Calibration & $\tfrac{1}{n}\sum_c (\hat{p}_c - y_c)^2$. Lower better. \\
Feature Importance & Ranking & Permutation (primary) and Gini importance. \\
\bottomrule
\end{tabular}
\end{table}

\paragraph{AUC-ROC} primary; threshold-independent and robust to class
imbalance. \paragraph{F1} threshold-dependent; threshold tuned on
validation, applied unchanged to test. \paragraph{Brier} distinguishes
ranking accuracy from probability reliability.
\paragraph{Permutation importance} reported as primary because invariant
to scale and cardinality.

Three feature sets compared (Section~\ref{sec:ml-feature-eng}):
\begin{enumerate}
  \item \textbf{PM-only}: Features 1--8.
  \item \textbf{MSR-only}: Features 9--15.
  \item \textbf{Combined}: Features 1--15.
\end{enumerate}

A statistically meaningful improvement of PM-only or Combined over
MSR-only answers RQ3 affirmatively. Wilcoxon signed-rank applied across
48 repositories; effect size via Cliff's~$\delta$ following
Kampenes et al.~\cite{kampenes2007systematic}.


% =============================================================================
% 5.5 BASELINES
% =============================================================================
\section{Baseline Methods}
\label{sec:baselines}

\begin{table}[htbp]
\centering
\caption{Baseline methods.}
\label{tab:baselines}
\small
\begin{tabular}{lp{4cm}p{5cm}}
\toprule
\textbf{Baseline} & \textbf{Description} & \textbf{Rationale} \\
\midrule
B1: Historical Mean         & Mean lead time of completed cases in $\mathcal{L}_{\text{train}}$. & Naïve baseline. \\
\midrule
B2: Throughput MC           & Throughput-based MC~\cite{vacanti2015,savage2009flaw,magennis2015forecasting}. & State-of-practice; black-box. \\
\midrule
B3: Markov Analytical       & $\mathbb{E}[\text{LT}] = \sum_j N_{ij} \cdot \bar{d}_j$. & Process-aware but deterministic. \\
\midrule
B4: Historical Distribution & Empirical CDF resampling. & Distributional baseline. \\
\bottomrule
\end{tabular}
\end{table}

The baselines isolate (i)~process awareness (vs.\ B2),
(ii)~distributional vs.\ point (vs.\ B3), (iii)~stochastic model vs.\
naïve resampling (vs.\ B4).


% =============================================================================
% 5.6 STATISTICAL ANALYSIS
% =============================================================================
\section{Statistical Analysis Procedures}
\label{sec:statistical-analysis}

\subsection{Within-Repository Comparison}
\label{sec:within-repo}

Paired tests on the same test cases.

\textbf{Normality assessment:} Shapiro–Wilk on per-case error differences.
If normal ($p > 0.05$), paired $t$-test; else Wilcoxon signed-rank.

\textbf{Effect size:} Cohen's $d$ (parametric) or Cliff's $\delta$
(non-parametric)~\cite{kampenes2007systematic}.


\subsection{Cross-Repository Aggregation}
\label{sec:cross-repo-stats}

\begin{enumerate}
    \item Per-repository metric (e.g., CRPS) for proposed and each
    baseline.
    \item Treat 48 metric values as paired samples.
    \item Wilcoxon signed-rank across repositories.
    \item Report median improvement and IQR.
\end{enumerate}

\textbf{Multiple comparison correction:} Bonferroni,
$\alpha_{\text{adj}} = 0.05 / 4 = 0.0125$.


\subsection{Statistical Power Analysis}
\label{sec:power-analysis}

A priori power following
Kampenes et al.~\cite{kampenes2007systematic}.

\textbf{RQ2/RQ3 — Cross-repository ($n=48$).} Under
$\alpha_{\text{adj}} = 0.0125$, power for medium effect ($r = 0.3$) is
$1-\beta \approx 0.72$; large ($r = 0.5$): $> 0.97$. Effect sizes are
informed by Rogge-Solti and
Weske~\cite{RoggeSolti2015NonMarkovianSPN}. Power computed via
\begin{equation}
\label{eq:power}
1 - \beta = P(Z > z_{\alpha/2} - r\sqrt{n-1})
    + P(Z < -z_{\alpha/2} - r\sqrt{n-1})
\end{equation}
with $z_{0.00625} \approx 2.50$.

\begin{table}[htbp]
\centering
\caption{A priori power estimates ($n = 48$, four baselines).}
\label{tab:power-analysis}
\small
\begin{tabular}{lccc}
\toprule
\textbf{Effect size} & \textbf{$r$} &
  \textbf{Power ($\alpha = 0.05$)} &
  \textbf{Power ($\alpha_{\text{adj}} = 0.0125$)} \\
\midrule
Small         & 0.1 & 0.17 & 0.09 \\
Small--medium & 0.2 & 0.44 & 0.27 \\
Medium        & 0.3 & 0.75 & 0.55 \\
Medium--large & 0.4 & 0.92 & 0.80 \\
Large         & 0.5 & 0.98 & 0.94 \\
\bottomrule
\end{tabular}
\end{table}

The dataset is adequately powered for medium-to-large effects at the
adjusted level. Low power for small effects ($r \leq 0.2$) is accepted
since hypotheses concern practically meaningful improvements.

\textbf{RQ3 — Within-repository ML.} Median expected test-set size
80--120 cases (227K commits, 1:4 issue ratio); sufficient for AUC-ROC
which stabilises above $n \approx 50$~\cite{hanley1982meaning}.


\subsection{Reporting Standards}
\label{sec:reporting}

Per Wohlin et al.~\cite{wohlin2012experimentation}:
\begin{itemize}
    \item Median, IQR, box plots (preferred over mean for skewed
    distributions).
    \item Test name, statistic, $p$-value, effect size for every
    comparison.
    \item 95\% CIs for all aggregate metrics.
    \item Calibration plots, CRPS box plots, feature importance rankings.
\end{itemize}


% =============================================================================
% 5.7 THREATS TO VALIDITY
% =============================================================================
\section{Threats to Validity}
\label{sec:threats}

Following Wohlin et al.~\cite{wohlin2012experimentation}.

\subsection{Internal Validity}
\label{sec:threats-internal}

\textbf{Temporal information leakage.} Mitigated by strict temporal
hold-out (Definition~\ref{def:temporal-holdout}).

\textbf{Confounding by project phase.} Partially mitigated by 70/30 split
plus sensitivity analysis with alternative ratios.

\textbf{Case correlation errors.} Mitigated by correlation quality
metrics (Definition~\ref{def:corr-quality}) and the cascade strategy.

\textbf{Bot contamination.} Mitigated by Stage~1 bot filtering plus
sensitivity analysis with/without filtering on a subset.


\subsection{External Validity}
\label{sec:threats-external}

\textbf{Open-source bias.} Results interpreted as feasibility evidence,
not absolute accuracy applicable to all settings.

\textbf{GitHub platform bias.} Stage~1 designed with pluggable extractors
for alternative platforms (GitLab, Bitbucket, Azure DevOps), but
extensibility not empirically validated here.

\textbf{Language and domain bias.} Partially addressed by cross-repository
generalization protocol (Section~\ref{sec:cross-repo}).


\subsection{Construct Validity}
\label{sec:threats-construct}

\textbf{Lead time definition.} First-event to last-event measurement.
Justified by data availability and consistency with PM-based measurement.

\textbf{Delay threshold for ML labels.} Three threshold strategies
evaluated (Section~\ref{sec:ml-labels}); results reported per strategy.

\textbf{Metric selection.} CRPS chosen as primary; MAE reported
alongside for comparability with point-forecast baselines.


\subsection{Conclusion Validity}
\label{sec:threats-conclusion}

\textbf{Statistical power.} Mitigated by:
(i)~minimum case count condition in Equation~\ref{eq:acceptance}
excluding $|\mathcal{L}_{\text{test}}| < 50$ from E2;
(ii)~cross-repository aggregation on $n = 48$ pairs achieving
$1-\beta \geq 0.72$ for medium effects at
$\alpha_{\text{adj}} = 0.0125$
(Section~\ref{sec:power-analysis}, Table~\ref{tab:power-analysis});
(iii)~per-repository CIs reported alongside aggregates.

\textbf{Multiple comparisons.} Bonferroni correction
(Section~\ref{sec:cross-repo-stats}).

\textbf{Reproducibility.} Released artifacts: (i)~repository identifiers
and dates; (ii)~PATHCAST source; (iii)~evaluation scripts;
(iv)~intermediate data (event logs, transition matrices, simulation
results). Raw repository data not redistributed but
re-collectable via GitHub API~\cite{github2023api} from
Appendix~\ref{app:repository-list}.


% =============================================================================
% 5.8 CHAPTER SUMMARY
% =============================================================================
\section{Chapter Summary}
\label{sec:methodology-summary}

\begin{table}[htbp]
\centering
\caption{Summary: research questions, evaluation stages, metrics, baselines.}
\label{tab:methodology-summary}
\small
\begin{tabular}{llll}
\toprule
\textbf{RQ} & \textbf{Stage} & \textbf{Primary Metrics} & \textbf{Baselines} \\
\midrule
RQ1 & E1 & Fitness, precision, coverage    & N/A (quality) \\
RQ2 & E2 & CRPS, ACE, MAE                  & B1, B2, B3, B4 \\
RQ3 & E2 & AUC-ROC, F1, Brier              & PM vs.\ MSR vs.\ Combined \\
\bottomrule
\end{tabular}
\end{table}

The evaluation answers three questions: Can the pipeline produce usable
event logs and models from real repositories (E1, RQ1)? Do process-aware
stochastic forecasts outperform process-blind baselines (E2, RQ2)? Do
process-derived features add predictive value beyond traditional metrics
(E2, RQ3)?
